\section{Introduction and prior work}\label{sec:introduction}

An orbit-counting function may have the expected exponential order of
growth without having a single asymptotic constant. For a map $f$ with
finitely many fixed points of every iterate, let $\pi_f(N)$ be the number
of primitive periodic orbits of length at most the integer $N$. Given an
exponential growth base $\Lambda>1$, the observable studied here is the
set of accumulation points of
\begin{equation}\label{eq:normalized_count_intro}
  \Pi_f(N)=\frac{N\pi_f(N)}{\Lambda^N},\qquad N\in\Z_{\geq1}.
\end{equation}
Arithmetic factors in the fixed-point sequence can prevent this sequence
from converging. Describing all its subsequential limits is then a
different question from determining its order of magnitude.

Byszewski, Cornelissen and Houben introduced finite-adelically distorted
systems, abbreviated FAD systems, to treat fixed-point formulas with
finitely many valuation-dependent factors. Their detector theorem realizes
the accumulation set in~\eqref{eq:normalized_count_intro} as the image of
an explicit continuous function on a compact group
\cite[Definitions 7.1.1--7.1.2 and Theorem 12.4.3]{bch2024}.
Their word \emph{hyperbolic} means that the multiplicative part has a
unique dominant root. In that case the archimedean phase disappears,
leaving a finite residue coordinate and finitely many $p$-adic coordinates.
For one distortion prime and no wild exponential factor, they prove that
the image is a union of a finite set and a Cantor set. Problem 14.1.1 in
their public version of 19 April 2024 singles out the remaining hyperbolic
cases: at least two distortion primes, or a nonzero wild term
\cite[Theorem 12.5.2, Remark 12.5.3 and Problem 14.1.1]{bch2024}.

This article proves that every distorted hyperbolic FAD accumulation set
is itself a Cantor set and is metrically thinner than the topological
classification alone indicates. If $d\geq1$ distortion primes are active,
then
\begin{equation}\label{eq:cover_intro}
  \neps(\LL_f)\leq C\bigl(1+\log(1/\eps)\bigr)^{2d}
  \qquad(0<\eps<\eps_0),
\end{equation}
where $\neps$ counts intervals of length $\eps$. In particular,
$\boxdim\LL_f=\hausdim\LL_f=0$.
The constants are allowed to depend on the fixed system. No uniformity
over arbitrarily large prime sets or exponents approaching zero is asserted.
Theorem~\ref{thm:image} proves a standalone image theorem under slightly
weaker periodicity hypotheses than the FAD definition.
Corollary~\ref{cor:fad} supplies the dynamical conclusion.

The detector viewpoint already has close predecessors. For $S$-integer
dynamics, Everest, Miles, Stevens and Ward control normalized orbit counts
through compact-group rotations and obtain a quantitative one-prime
injectivity result \cite[Theorem 1.1, Lemma 2.4 and Corollary 2.5]{emsw2007}.
For abelian varieties in positive characteristic, in the
unique-dominant-root setting and the branch that is not very inseparable,
Byszewski and Cornelissen prove the one-characteristic finite-set/Cantor-set
conclusion through a one-variable detector and local injectivity
\cite[Proposition 9.4 and Theorem 9.5]{bc2018}.
These counting reductions and examples are classical inputs to the present
problem. The term hyperbolic in this article is always used in the precise
FAD sense of~\cite[Definition 10.3.9]{bch2024}, not inferred from the
terminology in the title of~\cite{emsw2007}.

\begin{table}[t]
\centering
\small
\begin{tabularx}{\textwidth}{@{}p{0.27\textwidth}XX@{}}
\toprule
Regime & BCH public v2 conclusion & Conclusion here\\
\midrule
Hyperbolic, no distortion & Finite accumulation set & Same finite image\\
Hyperbolic, one prime, no wild term & Finite set united with a Cantor set
  & Cantor set; bound~\eqref{eq:cover_intro}\\
Hyperbolic, multiprime or wild & Cardinality continuum; topology left open
  & Cantor set; bound~\eqref{eq:cover_intro}\\
Nonhyperbolic & Contains a nondegenerate interval & Not treated\\
\bottomrule
\end{tabularx}
\caption{Theorem scope, with the prior-work column referring specifically
to~\cite[Theorems 12.5.1--12.5.2 and Problem 14.1.1]{bch2024}.
The distorted rows assume at least one active exponent. The covering
estimate is an upper bound, not an asserted sharp scale.}
\label{tab:scope}
\end{table}

Two independent arguments are needed. A continuous image of a profinite
space need not remain totally disconnected, and even a zero-dimensional
compact set may have isolated points. The covering proof retains only the
first $L$ summands of the detector series and refines a valuation tree only
along balls containing one of their centers. At depth $K$ this needs
$O(LK)$ pieces, instead of all $p^K$ residue classes. Taking both $L$ and
$K$ logarithmic in $1/\eps$ proves~\eqref{eq:cover_intro}.

Perfectness comes from a positive radial-kernel lemma. Each active kernel
has strictly negative Fourier coefficients at every nontrivial conductor.
Among finitely many exponent types, one has the slowest coefficient decay.
High-conductor characters approach one at each fixed ordinary integer,
and summability then prevents cancellation among positive translated
kernels. Local rescaling gives nonconstancy in every relevant $p$-adic
ball. To apply this argument in several coordinates, the other coordinates
are fixed at negative integers. That selection device already occurs in
BCH's proof of their cardinality dichotomy
\cite[proof of Theorem 12.5.1, pp.~132--133]{bch2024}.
The additional conclusion here follows from applying the Fourier lemma
on \emph{every} detector cylinder and using the coprime exponent periods
to supply an active center there.

The source comparison is deliberately version-specific. Cornelissen and
Park study the Ces\`aro mean of normalized fixed-point counts and large
deviations for orbit decompositions \cite[Theorem C and Section 4]{cp2026}.
Their reference 6 lists the BCH work as forthcoming at EMS Press in 2026.
The final book text was not available in the source audit for this article.
Thus the statement that the present proof addresses the open hyperbolic
regimes refers to the checked 2024 public version; it is not an unconditional
priority claim over subsequent book revisions. This qualification is
separate from the mathematical statement proved below.

Section~\ref{sec:image} states the image theorem.
Sections~\ref{sec:covers}--\ref{sec:perfectness} prove its metric and
topological parts. Section~\ref{sec:dynamics} applies it to genuine FAD
systems, including multiprime and wild examples. The final section records
the boundaries of the result.
